
\documentclass[11pt,a4paper]{article}

% ------------------------------------------------
% Page layout
% ------------------------------------------------
\usepackage[a4paper,margin=2.5cm]{geometry}

% ------------------------------------------------
% Mathematics and lists
% ------------------------------------------------
\usepackage{amsmath,amssymb}
\usepackage{enumitem}

% ------------------------------------------------
% URLs and hyperlinks
% ------------------------------------------------
\usepackage{xurl}
\usepackage[hidelinks]{hyperref}

% ------------------------------------------------
% Persian typesetting
% ------------------------------------------------
\usepackage{xepersian}

% ------------------------------------------------
% Fonts
% AutoFakeBold simulates bold if a bold font face is unavailable.
% ------------------------------------------------
\settextfont[AutoFakeBold=2]{B Nazanin}
\setlatintextfont[AutoFakeBold=2]{Times New Roman}

% ------------------------------------------------
% Formatting
% ------------------------------------------------
\setlength{\parindent}{0pt}
\setlength{\parskip}{0.5em}
\setlength{\emergencystretch}{3em}
\linespread{1.25}

% ------------------------------------------------
% English text
% ------------------------------------------------
\newcommand{\EN}[1]{\lr{#1}}

% ------------------------------------------------
% Lists
% ------------------------------------------------
\setlist[itemize]{
	label=\textbullet,
	leftmargin=1.8em,
	itemsep=0.4em,
	topsep=0.4em,
	parsep=0pt
}

\setlist[enumerate]{
	label=\arabic*.,
	leftmargin=2em,
	itemsep=0.5em,
	topsep=0.4em,
	parsep=0pt
}

% ------------------------------------------------
% Document
% ------------------------------------------------
\begin{document}
	
	% =================================================
	% Title
	% =================================================
	
	\begin{center}
		
		{\large\bfseries پیشنهاد پروژه درس محاسبات تکاملی}\\[0.8em]
		
		{\Large\bfseries
			استراتژی‌های تکاملی
			\EN{(Evolution Strategies)}
			به‌عنوان جایگزینی مقیاس‌پذیر برای یادگیری تقویتی
		}\\[1em]
		
		{\normalsize
			\textbf{نماینده و اعضای گروه:}
			محمدامین نقیان و علیرضا مهدی‌زاده
		}
		
	\end{center}
	
	\vspace{0.4em}
	\hrule
	\vspace{0.7em}
	
	% =================================================
	% Main paper
	% =================================================
	
	\textbf{مقاله اصلی}
	
	\begin{latin}
		
		T. Salimans, J. Ho, X. Chen, S. Sidor, and I. Sutskever.
		
		\textit{Evolution Strategies as a Scalable Alternative to
			Reinforcement Learning}.
		
		OpenAI, 2017.
		
		\url{https://doi.org/10.48550/arXiv.1703.03864}
		
	\end{latin}
	
	\vspace{0.4em}
	\hrule
	\vspace{0.7em}
	
	% =================================================
	% Section 1
	% =================================================
	
	\section*{۱. تشریح مسئله، ایده الگوریتم و خروجی نهایی}
	
	در مسائل یادگیری تقویتی، هدف آموزش یک عامل، مانند عامل
	کنترل‌کننده حرکت ربات یا عامل بازی‌کننده، برای بیشینه‌سازی
	پاداش تجمعی در طول زمان است. با این حال، بهینه‌سازی سیاست
	در این مسائل می‌تواند با چالش‌هایی مانند پرنوسان بودن تخمین
	گرادیان و هزینه بالای تعامل با محیط مواجه شود.
	
	در مقاله مورد بررسی
	\EN{(Salimans et al., 2017)}،
	نشان داده می‌شود که می‌توان از یک روش بهینه‌سازی مبتنی بر
	جمعیت، به نام استراتژی‌های تکاملی
	\EN{(Evolution Strategies (ES))}،
	برای آموزش سیاست‌های کنترلی استفاده کرد.
	
	در این روش، به‌جای محاسبه مستقیم گرادیان از طریق پس‌انتشار
	خطا در طول زمان، نسخه‌های مختلفی از پارامترهای سیاست با
	افزودن نویز تصادفی تولید و در محیط ارزیابی می‌شوند. سپس،
	با استفاده از پاداش حاصل از این ارزیابی‌ها، تخمینی از گرادیان
	مورد انتظار تابع هدف به دست می‌آید و پارامترهای سیاست
	به‌روزرسانی می‌شوند.
	
	یکی از مزایای مهم این روش، امکان موازی‌سازی ارزیابی سیاست‌ها
	روی چندین هسته پردازنده است. البته کارایی آن به نوع محیط،
	بودجه محاسباتی و نحوه پیاده‌سازی نیز وابسته است.
	
	خروجی نهایی پروژه شامل یک پیاده‌سازی پایتونی از الگوریتم
	استراتژی‌های تکاملی و اجرای موازی آن روی هسته‌های پردازنده
	خواهد بود. این پیاده‌سازی در چند محیط یادگیری تقویتی ارزیابی
	می‌شود.
	
	نتایج در قالب نمودارهای مقایسه‌ای عملکرد، شامل میانگین پاداش
	برحسب تعداد گام‌های محیط و زمان اجرای واقعی، همراه با
	بازه‌های اطمینان آماری ارائه خواهند شد. همچنین، تأثیر برخی
	مؤلفه‌های کلیدی الگوریتم، مقیاس‌پذیری محاسباتی و عملکرد عامل
	آموزش‌دیده بررسی می‌شود. در نهایت، یک ویدئوی کوتاه از عملکرد
	عامل نیز به‌عنوان بخشی از خروجی پروژه ارائه خواهد شد.
	
	% =================================================
	% Section 2
	% =================================================
	
	\section*{۲. بنچمارک‌های پیشنهادی و خطوط پایه مقایسه‌ای}
	
	از آنجا که مقیاس آزمایش‌های مقاله اصلی، شامل ارزیابی روی
	بازی‌های آتاری و محیط‌های شبیه‌سازی فیزیکی با منابع محاسباتی
	گسترده، با امکانات سخت‌افزاری ما همخوانی ندارد، آزمایش‌ها
	ابتدا روی محیط‌های استاندارد و سبک‌تر کتابخانه
	\EN{Gymnasium} انجام می‌شوند.
	
	\begin{enumerate}
		
		\item \textbf{محیط‌های تست اولیه و اعتبارسنجی:}
		
		محیط‌های \EN{CartPole-v1} و \EN{LunarLander-v3}
		برای کنترل گسسته و محیط \EN{Pendulum-v1}
		برای کنترل پیوسته.
		
		\item \textbf{محیط‌های مقایسه اصلی:}
		
		محیط‌های شبیه‌سازی فیزیکی \EN{HalfCheetah-v5}
		و \EN{Hopper-v5}، مشروط به سازگاری نسخه‌های
		وابستگی و امکان اجرای آن‌ها روی سخت‌افزار موجود.
		
		\item \textbf{معیارهای ارزیابی:}
		
		میانگین پاداش سیاست آموزش‌دیده در ازای بودجه محاسباتی
		مشخص، تعداد گام‌های محیط و زمان اجرای واقعی
		\EN{(Wall-clock time)}.
		
	\end{enumerate}
	
	برای ارزیابی مزایای الگوریتم، دو خط پایه
	\EN{(Baseline)} به شرح زیر در نظر گرفته می‌شوند:
	
	\begin{enumerate}
		
		\item \textbf{خط پایه اول: تپه‌نوردی تصادفی
			\EN{(Random Hill Climbing)} یا \EN{(1+1)-ES}.}
		
		در هر تکرار، یک نسخه جهش‌یافته از پارامترهای سیاست
		تولید و ارزیابی می‌شود. اگر این نسخه عملکرد بهتری داشته
		باشد، جایگزین نسخه فعلی خواهد شد. این روش برای بررسی
		میزان سودمندی استفاده از جمعیت و تخمین گرادیان در روش
		اصلی مناسب است.
		
		\item \textbf{خط پایه دوم: \EN{REINFORCE}.}
		
		این روش به‌عنوان یکی از الگوریتم‌های پایه گرادیان سیاست
		در نظر گرفته می‌شود تا عملکرد استراتژی‌های تکاملی با یک
		روش متداول یادگیری تقویتی مقایسه شود. در صورت استفاده
		از خط پایه یا کاهش واریانس، جزئیات پیاده‌سازی آن به‌صورت
		شفاف گزارش خواهد شد.
		
	\end{enumerate}
	
	در صورت وجود فرصت کافی، نتایج با الگوریتم
	\EN{PPO} از کتابخانه
	\EN{Stable-Baselines3} نیز مقایسه خواهند شد.
	
	آزمایش‌ها با دست‌کم پنج بذر تصادفی مستقل
	\EN{(Random Seeds)} برای هر تنظیم اجرا می‌شوند.
	میانگین و پراکندگی نتایج گزارش خواهد شد و در صورت مناسب
	بودن فرض‌های آماری و طراحی آزمایش، آزمون ناپارامتری
	\EN{Mann--Whitney U} برای مقایسه نتایج به کار می‌رود.
	
\end{document}
